\chapter{算法基础与复杂性}

\section{领域导言与演进背景}

\begin{preamblebox}[知识领域概述与变革]
算法与数据结构是计算机科学与软件工程不可动摇的基石。任何理论计算模型与现实工程软件，其本质均是在具备特定物理与逻辑结构的数据元素上执行算法操作。面对复杂的现实工程挑战，选择并设计出兼具计算效率、资源节约（包括能耗与碳足迹）与社会公平性的算法，要求从业者不仅理解算法的经验表现，更要洞察其内在的理论极限与计算复杂性。这种理论素养能使计算机专业人员超越具体的编程语言、编程范式或特定底层硬件架构，直击计算问题的本质规律。

\textbf{自 CS2013 以来的核心变革：}
\begin{enumerate}
    \item \textbf{重命名与回归理论根基}：从 CS2013 的“算法与复杂性”拓展为“算法基础”，重新强调所有计算机专业毕业生必须牢固掌握计算理论基础（自动机、形式语言、可计算性与不可判定性）。
    \item \textbf{增设“算法与社会”核心单元}：面对算法在现代社会决策中的无所不在，专门增设该单元，系统涵盖算法偏见、公平性度量、差分隐私与绿色低能耗计算等前沿社会议题。
    \item \textbf{与软件开发基础（SDF）的协同映射}：明确划分了初级数据结构应用与高级算法分析的边界，形成循序渐进的进阶路径。
\end{enumerate}
\end{preamblebox}

\section{核心学时分配建议}

本知识领域在“向日葵模型”下的学时推荐如下表所示：

\begin{table}[H]
\centering
\setlength{\tabcolsep}{4pt}\fontsize{8pt}{9.6pt}\selectfont\begin{tabularx}{\columnwidth}{Y c c}
\toprule
\textbf{知识单元} & \textbf{计算机核心} & \textbf{领域核心} \\
\midrule
基础数据结构与算法 & 36 & 7 \\
算法策略与范式 & 4 & -- \\
计算复杂性分析 & 4 & 2 \\
计算模型与形式理论 & 9 & 23 \\
算法与社会伦理 & 4 & -- \\
\midrule
\textbf{合计} & \textbf{57 学时} & \textbf{32 学时} \\
\bottomrule
\end{tabularx}
\end{table}

\section{知识单元与详细大纲}

\subsection{基础数据结构与算法}
\textbf{【学时建议】}：36 计算机核心学时，7 领域核心学时。

本单元构建起所有计算机科学专业学生必须掌握的数据结构与核心算法底座，与 SDF（软件开发基础）和 MSF（离散数学）紧密交织。

\subsubsection{计算机核心主题}
\begin{itemize}
    \item \textbf{抽象数据类型 (ADT)}：包（Bag）、容器、字典（Dictionary）、列表（List）与集合（Set）的代数规范与行为不变量。
    \item \textbf{数组与线性表}：一维与多维数组存储映射；线性查找与二分查找算法及其边界条件。
    \item \textbf{链表结构}：单向链表、双向链表、循环链表及其指针操作安全性与内存局部性特征。
    \item \textbf{栈与队列}：后进先出（LIFO）与先进先出（FIFO）机制、双端队列、基于堆的优先队列。
    \item \textbf{哈希表与散列技术}：哈希函数设计原理、冲突消解策略（线性探测、二次探测、链地址法、一致性哈希初步与动态重哈希）。
    \item \textbf{字符串与文本处理}：高效模式匹配算法（KMP 算法、Boyer-Moore 算法）。
    \item \textbf{树形结构与遍历}：
    \begin{itemize}
        \item 二叉树、多叉树、二叉搜索树（BST）及其前序、中序、后序与层序遍历。
        \item 平衡二叉树与多路平衡树原理：AVL 树、红黑树、2-3 树、B 树/B+ 树的核心旋转与重平衡机制。
        \item 堆结构（最大堆/最小堆）与其在优先队列中的工程实现。
        \item 编码与贪心应用：霍夫曼树与霍夫曼编码（Huffman Coding）。
    \end{itemize}
    \item \textbf{经典排序算法族}：
    \begin{itemize}
        \item $O(n^2)$ 经典排序：选择排序、插入排序、冒泡排序及其稳定性和原地性分析。
        \item $O(n \log n)$ 高效排序：快速排序（基准元划分策略与最坏情况退化规避）、归并排序、堆排序。
        \item 线性时间排序：计数排序、桶排序、基数排序及其应用边界。
    \end{itemize}
    \item \textbf{图结构与图遍历算法}：
    \begin{itemize}
        \item 图的拓扑属性：有向/无向、有环/无环（DAG）、连通性、带权图。
        \item 图的机器表示：邻接矩阵与邻接表的时空权衡。
        \item 基础遍历：广度优先搜索（BFS）及其在无权最短路中的应用；深度优先搜索（DFS）及其在环检测、连通分量与拓扑排序中的应用。
        \item 最小生成树（MST）：Prim 算法与 Kruskal 算法（并查集 Union-Find 的应用）。
        \item 最短路径算法族：单源最短路径 Dijkstra 算法（贪心策略）、Bellman-Ford 算法（负权边检测）；全源最短路径 Floyd-Warshall 算法。
        \item 传递闭包与 Warshall 算法。
    \end{itemize}
    \item \textbf{现代核心交叉技术}：
    \begin{itemize}
        \item 密码学基础散列与非对称基础（SHA-256、RSA 的数论支撑与实现逻辑）。
        \item 差分隐私（Differential Privacy）在聚合统计算法中的注入机制。
        \item 算法不变量的形式化证明与前置/后置条件。
    \end{itemize}
\end{itemize}

\subsubsection{领域核心与选修主题}
\begin{itemize}
    \item \textbf{领域核心}：
    \begin{itemize}
        \item 分布式共识算法（Consensus Algorithms，如 Paxos、Raft、区块链工作量证明 PoW 与权益证明 PoS 机制）。
        \item 计算几何基础（凸包 Graham 扫描法、最近点对分治求解）。
        \item 运筹与线性规划（Simplex 单纯形法与对偶原理）。
        \item 近似算法（针对背包问题、旅行商 TSP 的常数比例近似）。
        \item 随机化算法（MaxCut 局部搜索、球与箱模型 Balls and Bins、跳跃表 SkipList）。
    \end{itemize}
    \item \textbf{拓展选修}：
    \begin{itemize}
        \item 量子计算算法（Deutsch-Jozsa、Shor 大数质因数分解、Grover 数据库量子搜索）。
        \item 信号处理算法（快速傅里叶变换 FFT、离散卷积计算）。
        \item 遗传算法与演化计算（Genetic Algorithms、群体智能优化）。
    \end{itemize}
\end{itemize}

\subsubsection{示范性学习成果与能力目标}
\begin{enumerate}
    \item \textbf{[解释]} 针对本单元列出的所有基础 ADT 与数据结构，精准阐述其数学定义、不变量、多维底层存储表征及支持的增删查改原子操作。
    \item \textbf{[应用]} 给定具体的工业级应用场景（如高频交易订单撮合、海量日志检索），选择并实现最优的数据结构，证明其时空渐进复杂度的优越性。
    \item \textbf{[评估]} 全面对比各种排序、图搜索算法在最坏、最好、平均情况下的时间开销与内存局部性，评估真实硬件缓存环境对理论性能的影响。
    \item \textbf{[开发]} 针对复杂现实需求（例如基于多因子约束的路径规划系统），组合设计出混合数据结构，并撰写工程可行性与健壮性论证报告。
\end{enumerate}

\subsection{算法策略与范式}
\textbf{【学时建议】}：4 计算机核心学时。

本单元系统提炼通用的算法设计模式与高级思维方法：
\begin{itemize}
    \item \textbf{穷举与蛮力法}：组合搜索空间生成与状态遍历。
    \item \textbf{分治法}：大问题独立子问题解构、合并技术（典型范例：二分查找、快速排序、归并排序、Strassen 矩阵乘法）。
    \item \textbf{减治法}：每次将问题规模递减常数或常数因子（典型范例：插入排序、二叉查找、拓扑排序）。
    \item \textbf{动态规划}：最优子结构与重叠子问题、备忘录自顶向下法与制表自底向上法（典型范例：最长公共子序列 LCS、背包问题、Floyd 算法）。
    \item \textbf{贪心策略}：贪心选择性质与无后效性证明（典型范例：哈夫曼编码、Dijkstra 最短路、Kruskal 最小生成树）。
    \item \textbf{回溯法与分支限界法}：状态剪枝、Alpha-Beta 剪枝与组合优化边界估算。
    \item \textbf{时空权衡与化归策略 (Space-Time Tradeoff \& Reduction)}：以空间换时间（散列表、预处理排序）与问题规约转化（2-3树转换、图问题向网络流转化）。
\end{itemize}

\subsection{计算复杂性分析}
\textbf{【学时建议】}：4 计算机核心学时，2 领域核心学时。

本单元建立算法效能的严谨数理评估框架：
\begin{itemize}
    \item \textbf{渐进记号体系}：$O$（渐进上界）、$\Omega$（渐进下界）、$\Theta$（紧确界）、$o$ 与 $\omega$ 的严谨数学定义及极限证明。
    \item \textbf{基础时间复杂度阶}：常数阶 $O(1)$、对数阶 $O(\log n)$、线性阶 $O(n)$、线性对数阶 $O(n \log n)$、多项式阶 $O(n^k)$、指数阶 $O(2^n)$、阶乘阶 $O(n!)$。
    \item \textbf{递推关系式求解}：主定理三大分支条件推导与应用、代入归纳法、递归树分析法。
    \item \textbf{均摊分析}：聚合分析、记账法、势能法（Potential Method，如动态扩容数组分析）。
    \item \textbf{难解性与计算复杂性理论}：
    \begin{itemize}
        \item 判定问题与搜索问题的多项式时间可解性（类 P）。
        \item 非确定性多项式时间（类 NP）与多项式时间验证机概念。
        \item NP-完全性与 Cook-Levin 定理、经典 21 个 NPC 问题规约链条（SAT、3-SAT、团问题、顶点覆盖、哈密顿回路、TSP、子集和）。
        \item P 对 NP 问题的哲学意义与计算局限。
    \end{itemize}
\end{itemize}

\subsection{计算模型与形式理论}
\textbf{【学时建议】}：9 计算机核心学时，23 领域核心学时。

回归计算理论核心，使学生掌握计算能力的本质边界：
\begin{itemize}
    \item \textbf{有穷自动机与正则语言}：确定性有穷自动机 (DFA)、非确定性有穷自动机 (NFA)、子集构造法等价转化、正则表达式与正则文法、泵引理 及其不可正则性证明。
    \item \textbf{上下文无关文法与下推自动机}：上下文无关文法产生式、语法分析树与歧义性、下推自动机 (PDA) 的栈式存储模型、非上下文无关语言的泵引理。
    \item \textbf{图灵机与丘奇-图灵论题 (Turing Machines \& Church-Turing Thesis)}：标准单带图灵机形式化定义、多带/非确定性图灵机等价性、通用图灵机、丘奇-图灵论题的深刻内涵。
    \item \textbf{可计算性与不可判定性 (Computability \& Undecidability)}：图灵可识别与图灵可判定语言、对角线法证明不可数性、停机问题 的不可判定性严密反证法、波斯特对应问题 (PCP) 与莱斯定理 (Rice's Theorem)。
\end{itemize}

\subsection{算法与社会伦理}
\textbf{【学时建议】}：4 计算机核心学时。

\begin{itemize}
    \item \textbf{算法偏见与歧视}：训练数据历史偏见与算法反馈循环；招聘筛选、司法量刑、信贷审批算法中的不公诱因与防范机制。
    \item \textbf{算法公平性度量标准}：人口统计学一致性、均等化赔率、预测均等 及其内在的数学互斥性定理。
    \item \textbf{透明度、可解释性与知情权 (Explainability \& Right to Explanation)}：黑箱算法决策的审计追踪、模型决策证据生成机制。
    \item \textbf{绿色计算与算法能耗}：大型模型训练与高并发图计算的碳排放评估、算力节约型与能耗感知型算法设计。
\end{itemize}

\section{专业素养与思维倾向}

在算法领域的学习与实践中，重点塑造学生的以下专业特质：
\begin{itemize}
    \item \textbf{严密分析性}：面对模糊混沌的现实需求，能剥离表象，建立精确的形式化数学模型。
    \item \textbf{精益求精与严谨性}：深入推敲极端边界情况、死循环风险与内存开销溢出。
    \item \textbf{坚韧不拔的探索精神}：在攻坚复杂 NP-难问题或多约束系统时，能持续尝试近似、启发式或松弛技术。
    \item \textbf{负责任的工程伦理}：时刻警惕算法可能引发的社会排斥、隐私泄露与能耗过载风险。
\end{itemize}

\section{数学基础要求}

算法领域对离散数学与应用数学有着高度依赖：
\begin{itemize}
    \item \textbf{必须具备的离散数学核心}：集合论、谓词逻辑、二元关系（等价关系、偏序关系）、图论（度数、回路、树的同构）、初等组合计数、形式化数学证明（直接反证法、良序原理、超限数学归纳法）。
    \item \textbf{概率论与数理统计}：期望与方差的线性性质、条件概率、大数定律、贝叶斯定理、马尔可夫不等式与切比雪夫不等式（用于随机算法与哈希性能分析）。
    \item \textbf{线性代数与微积分}：矩阵运算、特征值、函数的渐近极限比较（L'H\^opital 法则在复杂度比较中的应用）。
\end{itemize}

\section{课程设计与开课规划建议}

联合工作组推荐两套典型开课方案：

\subsection{三门课贯穿方案 (面向通用型计算机专业)}
\begin{enumerate}
    \item \textbf{CS1：计算机编程基础 I}（涵盖基础数组、字符串查找、基础对象与迭代遍历，渗透初阶时空直觉）。
    \item \textbf{CS2：面向对象与数据结构}（15 学时覆盖单/双链表、栈、队列、哈希表、二叉平衡树、快速排序、归并排序、SHA-256 算法及 2 学时算法社会学）。
    \item \textbf{CS3：算法分析与计算理论综合课 (Algorithms \& Computation)}（以 16 学时高阶图论、动态规划、NPC 理论结合 9 学时计算模型与自动机基础）。
\end{enumerate}

\subsection{四门课深化方案 (面向研究型与深层次理论方向)}
将理论计算模型单独开设为一门 32 学时完整的《计算理论形式化方法》（涵盖自动机、下推系统、图灵机、不可判定性与复杂性类深入），而在《高级算法设计与分析》中腾出学时专攻均摊分析、网络流、近似算法与几何算法。

\section{真实能力任务规范}

\begin{taskbox}[任务 1：急救车辆智能紧急调度路由系统]
\textbf{【任务情境】}：设计一个在复杂城市路网拥堵与突发施工事件下的救护车智能调度路由模块，以毫秒级响应将危重病人送往最优急救医院。\\
\textbf{【能力陈述】}：检索、评估并融合现存高效图算法（如动态加权 Dijkstra 算法、A* 启发式导航算法结合收缩层次结构 Contraction Hierarchies），解决带实时权值变动的计算问题。\\
\textbf{【能力领域】}：软件工程 \quad \textbf{【能力单元】}：系统设计与评估\\
\textbf{【核心要求】}：计算机核心 \quad \textbf{【达标等级】}：\textbf{评估}
\end{taskbox}

\begin{taskbox}[任务 2：高精度慈善募捐个性化推荐与资产配置]
\textbf{【任务情境】}：根据用户输入的多元化价值偏好与税务属性，从数千家合规慈善机构中自动权衡并推荐最佳组合。\\
\textbf{【能力陈述】}：综合运用多目标贪心策略、带权二分图最大权匹配与多维背包动态规划模型，设计并实现一套高吞吐、抗并发的算法求解器。\\
\textbf{【能力领域】}：软件工程 \quad \textbf{【能力单元】}：系统开发\\
\textbf{【核心要求】}：计算机核心 \quad \textbf{【达标等级】}：\textbf{开发}
\end{taskbox}

\begin{taskbox}[任务 3：遗留客户服务系统的算法性能重构]
\textbf{【任务情境】}：某企业老旧工单处理系统在千万级数据量下出现严重查询阻塞，需彻底消除 $O(n^2)$ 级隐藏瓶颈。\\
\textbf{【能力陈述】}：剖析系统瓶颈来源，使用红黑树平衡索引、布隆过滤器与缓存哈希表替换原生遍历，将核心响应时间降低数个数量级。\\
\textbf{【能力领域】}：软件工程 \quad \textbf{【能力单元】}：系统重构与性能改进\\
\textbf{【核心要求】}：计算机核心 \quad \textbf{【达标等级】}：\textbf{开发}
\end{taskbox}

\begin{taskbox}[任务 4：跨机构医疗患者全流程状态跟踪自动化]
\textbf{【任务情境】}：构建一套覆盖门诊预检、手术排期、病房流转直至出院随访的全生命周期严密跟踪算法。\\
\textbf{【能力陈述】}：基于确定性有穷状态机 (DFA) 与层次状态图建立无死锁、无异常旁路的业务状态转移模型，并做全分支可达性形式化验证。\\
\textbf{【能力领域】}：软件工程 \quad \textbf{【能力单元】}：需求分析与架构设计\\
\textbf{【核心要求】}：计算机核心 \quad \textbf{【达标等级】}：\textbf{开发}
\end{taskbox}

\begin{taskbox}[任务 5：多维度复杂排班优化问题的难解性论证]
\textbf{【任务情境】}：客户提出一个兼顾多班次工时、请假优先级、技能准入与奖惩成本的万能实时排班需求。\\
\textbf{【能力陈述】}：撰写专业技术白皮书，通过多项式时间规约证明该问题为强 NP-难（规约到 3-SAT 或精确覆盖问题），指出其精确最优求解在多项式时间内不可行，并推荐高效模拟退火或遗传近似求解方案。\\
\textbf{【能力领域】}：理论计算 \quad \textbf{【能力单元】}：可行性评估与技术论证\\
\textbf{【核心要求】}：计算机核心 \quad \textbf{【达标等级】}：\textbf{解释与应用}
\end{taskbox}

\begin{taskbox}[任务 6：数据最优无损压缩极限的理论不可判定性论证]
\textbf{【任务情境】}：某企业试图研发一种“对任意输入数据均能产生理论极限最优压缩率且保证无损还原”的单体通用算法。\\
\textbf{【能力陈述】}：撰写严谨的理论学术报告，通过反证法证明该最优性判定问题可归约为图灵机停机问题（柯尔莫哥洛夫复杂性不可计算定理），以严密逻辑澄清科学边界，避免企业投入无谓研发资本。\\
\textbf{【能力领域】}：理论计算 \quad \textbf{【能力单元】}：技术论证\\
\textbf{【核心要求】}：计算机核心 \quad \textbf{【达标等级】}：\textbf{评估}
\end{taskbox}

\begin{taskbox}[任务 7：全民点赞与建议评级推荐算法的社会与伦理审计]
\textbf{【任务情境】}：社区自治平台拟推行基于热度点赞的市民建议自动排序与采纳机制。\\
\textbf{【能力陈述】}：撰写算法社会影响评估报告，揭示马太效应与沉默螺旋效应，设计差分隐私噪声注入与长尾代表性补偿模型，防止少数群体利益被算法遮蔽。\\
\textbf{【能力领域】}：软件与社会 \quad \textbf{【能力单元】}：社会议题调适\\
\textbf{【核心要求】}：计算机核心 \quad \textbf{【达标等级】}：\textbf{解释与评估}
\end{taskbox}

\begin{strip}
\section{附录：核心主题与技能要求等级对照表}
\noindent\begin{tabularx}{\textwidth}{p{0.18\textwidth} Y c c c}
\toprule
\textbf{知识单元} & \textbf{核心知识点与具体主题} & \textbf{认知等级} & \textbf{分类} & \textbf{学时} \\
\midrule
\multirow{4}{*}{\makecell[l]{算法/基础数据结构\\与算法}} 
& 常用 ADT（字典、映射、集合）、基础线性表与树 & 解释 & 计算机核心 & 12 \\
& 经典查找、快速排序、归并排序、堆排序实现与优化 & 应用 & 计算机核心 & 12 \\
& 图的连通性、Dijkstra 最短路、Prim/Kruskal 最小生成树 & 应用 & 计算机核心 & 12 \\
& 线性规划单纯形法、凸包几何算法、近似与随机算法 & 评估 & 领域核心 & 7 \\
\midrule
\multirow{2}{*}{\makecell[l]{算法/算法策略与范式}} 
& 分治法、减治法、动态规划、贪心法、回溯与分支限界原理 & 解释 & 计算机核心 & 2 \\
& 针对工程问题选择或重构适宜的算法范式进行求解 & 应用 & 计算机核心 & 2 \\
\midrule
\multirow{3}{*}{\makecell[l]{算法/复杂性分析}} 
& 渐进符号体系（大 O、大 $\Omega$、大 $\Theta$）与主定理推导 & 解释 & 计算机核心 & 2 \\
& 递归与非递归算法的时空开销度量、经验评测 & 应用 & 计算机核心 & 2 \\
& P 与 NP 问题概念、NPC 经典问题规约链条、均摊分析 & 评估 & 领域核心 & 2 \\
\midrule
\multirow{3}{*}{\makecell[l]{算法/计算模型与\\形式理论}} 
& 正则语言、DFA 与 NFA 等价性、正则表达式转化 & 应用 & 计算机核心 & 5 \\
& 上下文无关文法、下推自动机、图灵机定义与停机问题 & 解释 & 计算机核心 & 4 \\
& 泵引理证明、非确定图灵机、莱斯定理、不可判定性层级 & 开发 & 领域核心 & 23 \\
\midrule
\multirow{2}{*}{\makecell[l]{算法/算法与社会}} 
& 算法偏见来源、公平性数学度量（统计均等 vs 均等化赔率） & 解释 & 计算机核心 & 2 \\
& 差分隐私扰动注入、绿色计算与算法能耗感知优化 & 评估 & 计算机核心 & 2 \\
\bottomrule
\end{tabularx}
\end{strip}
